% Part: turing-machines
% Chapter: machines-computations
% Section: introduction

\documentclass[../../../include/open-logic-section]{subfiles}

\begin{document}

\olfileid{tur}{mac}{int}
\olsection{પરિચય}

ઉદાહરણ તરીકે $\Nat$થી $\Nat$માં જતું કોઈ વિધેય \emph{સંગણનીય}
હોવાનો અર્થ શું? શરૂઆતના જવાબોમાંનો એક, અને સૌથી જાણીતો જવાબ,
એ છે કે કોઈ વિધેય ટ્યુરિંગ મશીન વડે ગણી શકાય તો તે સંગણનીય છે.
એલન ટ્યુરિંગે 1936માં આ કલ્પના રજૂ કરી. ટ્યુરિંગ મશીન
\emph{સંગણનના નિદર્શ}નું એક ઉદાહરણ છે---એ ``સંગણન પ્રક્રિયા''નો
વિચાર ગણિતીય રીતે ચોક્કસ વ્યાખ્યાયિત કરવાની રીત છે. આ વિચારનો
ચોક્કસ અર્થ શું છે તે ચર્ચાનો વિષય છે, પરંતુ ટ્યુરિંગ મશીનો
સંગણન પ્રક્રિયા નિર્દિષ્ટ કરવાની એક રીત છે એ બાબતે વ્યાપક સહમતિ
છે. ``ટ્યુરિંગ મશીન'' શબ્દથી ચાલતા ભાગો ધરાવતા ભૌતિક મશીનની
છબી ઊભી થાય છે, પરંતુ કડક અર્થમાં ટ્યુરિંગ મશીન શુદ્ધ ગણિતીય
રચના છે; તેથી તે સંગણન પ્રક્રિયાના વિચારનું આદર્શીકરણ કરે છે.
ઉદાહરણ તરીકે, આપણે ટ્યુરિંગ મશીનને લાગતા સમય કે સ્મૃતિની
જરૂરિયાત પર કોઈ મર્યાદા મૂકતા નથી: સંગણના માટે બ્રહ્માંડના
પરમાણુઓની સંખ્યા કરતાં વધુ સંગ્રહસ્થાન અથવા વધુ પગલાં જોઈએ
તોપણ ટ્યુરિંગ મશીન તે ગણતરી કરી શકે છે.

\begin{explain}
ટ્યુરિંગ મશીનને કોઈ વિશિષ્ટ પ્રકારના કલ્પિત ઉપકરણના પ્રોગ્રામ
તરીકે સમજવું કદાચ સૌથી યોગ્ય છે. આ ઉપકરણમાં એક \emph{ટેપ} અને
એક \emph{વાંચન-લેખન હેડ} હોય છે. ટ્યુરિંગ મશીનની આપણી આવૃત્તિમાં
ટેપ એક દિશામાં (જમણી તરફ) અનંત હોય છે અને \emph{ખાનાં}માં
વહેંચાયેલી હોય છે. દરેક ખાનામાં મર્યાદિત \emph{વર્ણમાળા}માંથી
એક પ્રતીક હોઈ શકે છે. આવી વર્ણમાળામાં સાન્ત સંખ્યાનાં ગમે તેટલાં જુદાં પ્રતીકો
હોઈ શકે, પરંતુ આપણે મુખ્યત્વે ત્રણથી કામ ચલાવીશું:
$\TMendtape$, $\TMblank$ અને $\TMstroke$. ઉપકરણ શરૂ થાય ત્યારે
ટેપ ખાલી હોય છે (એટલે દરેક ખાનામાં $\TMblank$ હોય છે), સિવાય કે
સૌથી ડાબા ખાનામાં $\TMendtape$ અને સાન્ત સંખ્યામાં અન્ય ખાનાંમાં
\emph{ઇનપુટ} હોય છે. કોઈ પણ સમયે ઉપકરણ સાન્ત સંખ્યાની
\emph{અવસ્થાઓ}માંથી એક અવસ્થામાં હોય છે. શરૂઆતમાં હેડ સૌથી
ડાબું ખાનું વાંચે છે અને ઉપકરણ નિર્દિષ્ટ \emph{પ્રારંભિક અવસ્થા}માં
હોય છે. ઉપકરણના ચાલવાના દરેક પગલે, હાલમાં હેડ જે ખાનું વાંચે
છે તેની સામગ્રી અને ઉપકરણની અવસ્થા મળીને, ટ્યુરિંગ મશીનના
પ્રોગ્રામ મુજબ આગળ શું થશે તે નક્કી કરે છે. ટ્યુરિંગ મશીનનો
પ્રોગ્રામ એક આંશિક વિધેય વડે અપાય છે: તે અવસ્થા~$q$ અને
પ્રતીક~$\sigma$ ઇનપુટ તરીકે લે છે અને ત્રિક
$\tuple{q', \sigma', D}$ આપે છે. ઉપકરણ અવસ્થા~$q$માં હોય અને
પ્રતીક~$\sigma$ વાંચે ત્યારે તે હાલના ખાનામાંનું પ્રતીક
$\sigma'$થી બદલે છે; $D$ અનુક્રમે $\TMleft$, $\TMright$ કે
$\TMstay$ હોય તે મુજબ હેડ ડાબે, જમણે કે એ જ સ્થાને જાય છે;
અને ઉપકરણ અવસ્થા~$q'$માં પ્રવેશે છે.

ઉદાહરણ તરીકે, \olref{fig:tm}માં દર્શાવેલી સ્થિતિ જુઓ.
\begin{figure}
  \olasset{\olpath/assets/diagrams/turing-machine.tikz}
  \caption{પોતાનો પ્રોગ્રામ ચલાવતું ટ્યુરિંગ મશીન.}
  \ollabel{fig:tm}
\end{figure}
ટ્યુરિંગ મશીનની ટેપનો દેખાતો ભાગ સૌથી ડાબા ખાનામાં ટેપના
છેડાનું પ્રતીક~$\TMendtape$ ધરાવે છે. ત્યાર પછી ત્રણ $1$,
એક $0$ અને બીજા ચાર $1$ છે. હેડ ડાબેથી ત્રીજું ખાનું વાંચે
છે, જેમાં $1$ છે, અને મશીન અવસ્થા~$q_1$માં છે---અપણે કહીએ
કે ``મશીન અવસ્થા~$q_1$માં $1$ વાંચે છે.'' જો ટ્યુરિંગ મશીનનો
પ્રોગ્રામ ઇનપુટ~$\tuple{q_1, 1}$ માટે ત્રિક
$\tuple{q_2, 0, \TMstay}$ આપે, તો મશીન હવે ત્રીજા ખાનાનો
$1$ બદલીને $0$ લખશે, વાંચન-લેખન હેડને એ જ સ્થાને રાખશે અને
અવસ્થા~$q_2$માં જશે. પછી જો પ્રોગ્રામ ઇનપુટ~$\tuple{q_2, 0}$
માટે $\tuple{q_3, 0, \TMright}$ આપે, તો મશીન હવે તે $0$ પર
ફરી $0$ લખશે (એટલે હેડની નીચેની ટેપની સામગ્રી હકીકતમાં
બદલાશે નહીં), એક ખાનું જમણે ખસશે અને અવસ્થા~$q_3$માં
પ્રવેશશે. આ રીતે આગળ ચાલે છે.

મશીન કોઈ અવસ્થા~$q_n$ અને પ્રતીક~$\sigma$ની એવી જોડી પાસે
પહોંચે કે જેના માટે $\tuple{q_n, \sigma}$ને અનુરૂપ કોઈ સૂચના
ન હોય---એટલે ઇનપુટ $\tuple{q_n,\sigma}$ પર સંક્રમણ વિધેય
અવ્યાખ્યાયિત હોય---
ત્યારે આપણે કહીએ છીએ કે મશીન \emph{થંભે} છે. બીજા શબ્દોમાં,
મશીન પાસે કરવા માટે કોઈ સૂચના રહેતી નથી અને તે બિંદુએ તેની
ક્રિયા બંધ થાય છે. થંભવું ક્યારેક વિશિષ્ટ થંભવાની
અવસ્થા~$h$થી દર્શાવાય છે. આ બાબત આગળ વધુ વિગતે બતાવીશું.
\end{explain}

\begin{digress}
ટ્યુરિંગના ``સંગણનીય સંખ્યાઓ અંગે'' (``On computable numbers'') લેખની
વિશેષતા એ છે કે તેમાં
માત્ર ઔપચારિક વ્યાખ્યા જ નથી, પરંતુ એ વ્યાખ્યા સંગણનીયતાની
સહજ કલ્પનાને પકડે છે એવી દલીલ પણ છે. વ્યાખ્યામાંથી એ સ્પષ્ટ
થવું જોઈએ કે ટ્યુરિંગ મશીન વડે સંગણનીય દરેક વિધેય સહજ અર્થમાં
સંગણનીય છે. ટ્યુરિંગે ઊલટી દિશા માટે---એટલે જેને આપણે
સ્વાભાવિક રીતે સંગણનીય ગણીએ એવું દરેક વિધેય આવા મશીન વડે
ગણી શકાય છે તે માટે---ત્રણ પ્રકારની દલીલો આપી. તેમના પોતાના
શબ્દોમાં તે આ પ્રમાણે છે:
\begin{enumerate}
\item સહજ સમજને સીધી અપીલ.
\item બે વ્યાખ્યાઓ સમકક્ષ છે તેની સાબિતી (નવી વ્યાખ્યા સહજ
  સમજને વધુ ગમે એવી હોય ત્યારે).
\item સંગણનીય સંખ્યાઓના મોટા વર્ગોનાં ઉદાહરણો આપવાં.
\end{enumerate}
આપણો હેતુ સંગણનીયતાને ``સૈદ્ધાંતિક રીતે'' વ્યાખ્યાયિત કરવાનો
છે, એટલે સમય અને સ્થાનની વ્યવહારુ મર્યાદાઓ ધ્યાનમાં લીધા
વિના. અલબત્ત, સંગણનીયતાની વ્યાપકતમ વ્યાખ્યા હાથવગી થયા
પછી મર્યાદિત સંસાધનો સાથેના સંગણનનો પણ અભ્યાસ કરી શકાય;
તે ``સંગણકીય જટિલતા'' તરીકે ઓળખાતા વિષયનું કેન્દ્ર છે.
\end{digress}

\begin{history}
એલન ટ્યુરિંગે 1936માં ટ્યુરિંગ મશીનોની શોધ કરી. તે સમયે તેમનો
રસ પ્રથમ ક્રમના તર્કશાસ્ત્રની નિર્ણેયતામાં હતો, છતાં તેમના
લેખને કમ્પ્યુટરની રચનાના પાયા અંગેનો નિર્ણાયક લેખ કહેવામાં
આવ્યો છે. તેમાં ટ્યુરિંગ સંગણનીય વાસ્તવિક સંખ્યાઓ પર ધ્યાન
આપે છે, એટલે એવી વાસ્તવિક સંખ્યાઓ કે જેના દશાંશ વિસ્તરણો
ગણી શકાય; પરંતુ તેઓ નોંધે છે કે પ્રાકૃતિક સંખ્યાઓ પરનાં
સંગણનીય વિધેયો વગેરે માટે પણ આ વિચારોને અનુકૂળ બનાવવાનું
અઘરું નથી. નોંધો કે પ્રથમ કાર્યરત સામાન્ય હેતુનો કમ્પ્યુટર
1941માં બનાવાયો તેના પૂરા પાંચ વર્ષ પહેલાંની આ વાત છે
(જર્મન કોનરાડ ઝૂઝેએ પોતાના માતાપિતાના ઘરના બેઠકખંડમાં તે
બનાવ્યો હતો); બ્લેચલી પાર્કમાં ટ્યુરિંગ અને તેમના સાથીઓએ
સંકેત ઉકેલતું કોલોસસ બનાવ્યું (1943) તેના સાત વર્ષ પહેલાં;
અમેરિકન ENIAC (1945)ના નવ વર્ષ પહેલાં; મૅન્ચેસ્ટરમાં પ્રથમ
બ્રિટિશ સામાન્ય હેતુનો કમ્પ્યુટર---Manchester Small-Scale
Experimental Machine---બન્યો (1948) તેના બાર વર્ષ પહેલાં;
અને અમેરિકનોએ BINAC (1949)નું પ્રથમ પરીક્ષણ કર્યું તેના તેર
વર્ષ પહેલાં. Manchester SSEM પ્રથમ સંગ્રહિત-પ્રોગ્રામવાળો
કમ્પ્યુટર હતો---અગાઉનાં મશીનોને દરેક નવા કામ માટે હાથેથી
ફરી તાર જોડવા પડતા હતા.
\end{history}

\end{document}
